\chapter{Implementazione SQL}

\section{Creazione delle tabelle}

\vspace{2em}
\textbf{UTENTI}

\lstinputlisting[
    language=SQL, 
    firstline=19, 
    lastline=31, 
    firstnumber=1, 
    caption={Creazione della tabella UTENTI.}
]{sql/creazione_db.sql}

\vspace{2em}
\textbf{CATEGORIE}

\lstinputlisting[
    language=SQL, 
    firstline=33, 
    lastline=39, 
    firstnumber=1, 
    caption={Creazione della tabella CATEGORIE.}
]{sql/creazione_db.sql}

\vspace{2em}
\textbf{INGREDIENTI}

\lstinputlisting[
    language=SQL, 
    firstline=41, 
    lastline=49, 
    firstnumber=1, 
    caption={Creazione della tabella INGREDIENTI.}
]{sql/creazione_db.sql}

\vspace{2em}
\textbf{PROMOZIONI}

\lstinputlisting[
    language=SQL, 
    firstline=51, 
    lastline=58, 
    firstnumber=1, 
    caption={Creazione della tabella PROMOZIONI.}
]{sql/creazione_db.sql}

\vspace{2em}
\textbf{RICETTE}

\lstinputlisting[
    language=SQL, 
    firstline=60, 
    lastline=81, 
    firstnumber=1, 
    caption={Creazione della tabella RICETTE.}
]{sql/creazione_db.sql}

\vspace{2em}
\textbf{ORDINI}

\lstinputlisting[
    language=SQL, 
    firstline=83, 
    lastline=94, 
    firstnumber=1, 
    caption={Creazione della tabella ORDINI.}
]{sql/creazione_db.sql}

\vspace{2em}
\textbf{CLASSIFICAZIONI}

\lstinputlisting[
    language=SQL, 
    firstline=96, 
    lastline=110, 
    firstnumber=1, 
    caption={Creazione della tabella CLASSIFICAZIONI.}
]{sql/creazione_db.sql}

\vspace{2em}
\textbf{DETTAGLI\_RICETTA}

\lstinputlisting[
    language=SQL, 
    firstline=112, 
    lastline=128, 
    firstnumber=1, 
    caption={Creazione della tabella DETTAGLI\_RICETTA.}
]{sql/creazione_db.sql}

\vspace{2em}
\textbf{DETTAGLI\_ORDINE}

\lstinputlisting[
    language=SQL, 
    firstline=130, 
    lastline=150, 
    firstnumber=1, 
    caption={Creazione della tabella DETTAGLI\_ORDINE.}
]{sql/creazione_db.sql}

\vspace{2em}
\textbf{SCONTI}

\lstinputlisting[
    language=SQL, 
    firstline=152, 
    lastline=171, 
    firstnumber=1, 
    caption={Creazione della tabella SCONTI.}
]{sql/creazione_db.sql}

\vspace{2em}
\textbf{RECENSIONI}

\lstinputlisting[
    language=SQL, 
    firstline=173, 
    lastline=191, 
    firstnumber=1, 
    caption={Creazione della tabella RECENSIONI.}
]{sql/creazione_db.sql}
\vspace{1em}

\section{Indici}

\lstinputlisting[
    language=SQL, 
    firstline=197, 
    lastline=208, 
    firstnumber=1, 
    caption={Indici principali dello schema.}
]{sql/creazione_db.sql}
\vspace{1em}

\section{Trigger}

\vspace{2em}
\textbf{Trigger per INSERT su SCONTI}

Impedisce l'inserimento se esiste già uno sconto con stessa promozione, stessa categoria e range di ingredienti sovrapposto con quello desiderato.

\lstinputlisting[
    language=SQL, 
    firstline=217, 
    lastline=234, 
    firstnumber=1, 
    caption={Trigger per INSERT su SCONTI.}
]{sql/creazione_db.sql}

\vspace{3em}
\textbf{Trigger per UPDATE su SCONTI}

Impedisce l'aggiornamento se esiste già uno sconto con stessa promozione, stessa categoria e range di ingredienti sovrapposto con quello desiderato.

\lstinputlisting[
    language=SQL, 
    firstline=239, 
    lastline=258, 
    firstnumber=1, 
    caption={Trigger per UPDATE su SCONTI.}
]{sql/creazione_db.sql}

\vspace{3em}
\textbf{Trigger per INSERT su RECENSIONI}

Ricalcola e aggiorna la MediaRecensioni della ricetta appena recensita.

\lstinputlisting[
    language=SQL, 
    firstline=263, 
    lastline=274, 
    firstnumber=1, 
    caption={Trigger per INSERT su RECENSIONI.}
]{sql/creazione_db.sql}

\vspace{3em}
\textbf{Trigger per DELETE su RECENSIONI}

Ricalcola e aggiorna la MediaRecensioni della ricetta a cui la recensione eliminata faceva riferimento.

\lstinputlisting[
    language=SQL, 
    firstline=279, 
    lastline=289, 
    firstnumber=1, 
    caption={Trigger per DELETE su RECENSIONI.}
]{sql/creazione_db.sql}

\vspace{3em}
\textbf{Trigger per UPDATE su INGREDIENTI}

Ricalcola e aggiorna il prezzo di tutte le ricette che contengono l'ingrediente aggiornato (solo se è cambiato il prezzo di quest'ultimo).

\lstinputlisting[
    language=SQL, 
    firstline=294, 
    lastline=312, 
    firstnumber=1, 
    caption={Trigger per UPDATE su INGREDIENTI.}
]{sql/creazione_db.sql}

\vspace{3em}
\textbf{Trigger per INSERT su DETTAGLI\_RICETTA}

Ricalcola e aggiorna il NumeroIngredienti e il Prezzo della ricetta corrispondente.

\lstinputlisting[
    language=SQL, 
    firstline=318, 
    lastline=335, 
    firstnumber=1, 
    caption={Trigger per INSERT su DETTAGLI\_RICETTA.}
]{sql/creazione_db.sql}

\vspace{3em}
\textbf{Trigger per INSERT su ORDINI}

Impedisce l'inserimento se l'utente corrispondente non ha un IndirizzoSpedizione associato.

\lstinputlisting[
    language=SQL, 
    firstline=340, 
    lastline=354, 
    firstnumber=1, 
    caption={Trigger per INSERT su ORDINI.}
]{sql/creazione_db.sql}
\vspace{1em}

\section{Viste}

\vspace{1em}
\textbf{Fatturato giornaliero}

Per ogni giorno in cui sono stati effettuati ordini, mostra numero di ordini, quantita' totale venduta e incasso totale (tenendo conto degli sconti effettivamente applicati in fase d'ordine

\lstinputlisting[
    language=SQL, 
    firstline=365, 
    lastline=376, 
    firstnumber=1, 
    caption={Vista per il fatturato giornaliero.}
]{sql/creazione_db.sql}

\vspace{2em}
\textbf{Fatturato per ricetta}

Per ogni ricetta, quantita' totale ordinata, numero di ordini distinti in cui compare e incasso generato.
Utile per capire quali ricette "trainano" le vendite.

\lstinputlisting[
    language=SQL, 
    firstline=382, 
    lastline=394, 
    firstnumber=1, 
    caption={Vista per il fatturato per ricetta.}
]{sql/creazione_db.sql}

\vspace{2em}
\textbf{Recensioni negative}

Recensioni negative recenti (voto <= 3) delle ultime 2 settimane: utile per moderazione rapida o per contattare l'autore della ricetta.

\lstinputlisting[
    language=SQL, 
    firstline=400, 
    lastline=415, 
    firstnumber=1, 
    caption={Vista per le recensioni negative.}
]{sql/creazione_db.sql}